% Capítulo de soluções numéricas de EDOs (arquivos usados neste capítulo)
\def\arqEuler{MetodoEuler.ggb}
\def\arqRKdois{RungeKutta2.ggb}
\def\arqRKquatro{RungeKutta4.ggb}
\def\arqOrdem{OrdemMetodosEDO.ggb}
\def\arqSistemas{SistemasEDO_RK4.ggb}
\def\arqTiro{MetodoTiro.ggb}
\def\arqDifFin{DiferencasFinitasPVC.ggb}
\chapter[Soluções numéricas de EDOs]{Soluções numéricas de equações diferenciais ordinárias}\label{cap:edo}

\begin{tcolorbox}[colback=gray!4,colframe=gray!50,title={Como usar os aplicativos deste capítulo},fonttitle=\bfseries]
Ao longo do capítulo, os quadros laranja \emph{Atividade no GeoGebra} trazem um link
\textbf{[abrir o aplicativo]} que abre o arquivo correspondente no GeoGebra Classic:
\begin{itemize}[nosep]
  \item \ggbabrir{\arqEuler}{Método de Euler};
  \item \ggbabrir{\arqRKdois}{Runge-Kutta de 2ª ordem (Heun e ponto médio)};
  \item \ggbabrir{\arqRKquatro}{Runge-Kutta de 4ª ordem};
  \item \ggbabrir{\arqOrdem}{Ordem dos métodos: Euler, Heun e RK4};
  \item \ggbabrir{\arqSistemas}{Sistemas de EDOs com RK4};
  \item \ggbabrir{\arqTiro}{Método do tiro};
  \item \ggbabrir{\arqDifFin}{Diferenças finitas para problemas de contorno}.
\end{itemize}
Para que os links funcionem, mantenha este PDF e a pasta \texttt{ggb} juntos, na mesma
pasta, e abra o PDF num leitor que permita abrir arquivos externos (Adobe Acrobat Reader,
Okular, SumatraPDF, Foxit). Os leitores de PDF embutidos nos navegadores costumam bloquear
esse tipo de link; nesse caso, abra os arquivos da pasta \texttt{ggb} diretamente no GeoGebra.

Os quatro primeiros aplicativos começam com o problema $y'=y-x^2+1$, $y(0)=0{,}5$, que é o
exemplo desenvolvido no texto. Em todos é possível digitar outra equação, mudar a condição
inicial, o passo $h$ e o número de passos $n$, mostrar o campo de direções e usar o botão
\emph{Enquadrar o gráfico}. A janela da direita mostra as contas e tem barras de rolagem
(vertical e horizontal). A curva verde de referência é calculada pelo próprio aplicativo com
RK4 e um passo 40 vezes menor, o que basta para medir o erro dos métodos estudados.
\end{tcolorbox}

% ==================================================================
\section{Problemas de valor inicial}

Uma \emph{equação diferencial ordinária} (EDO) de primeira ordem relaciona uma função
desconhecida $y(x)$ com a sua derivada: $y'=f(x,y)$. Ela descreve como uma grandeza
varia: o crescimento de uma população, o resfriamento de um corpo, a carga de um capacitor, a
concentração de um medicamento no sangue. Uma EDO tem infinitas soluções; para escolher uma,
fixamos o valor de $y$ num ponto.

\begin{definicao}
Um \emph{problema de valor inicial} (PVI) de primeira ordem é
\[
y'=f(x,y),\qquad y(x_0)=y_0,
\]
e uma solução no intervalo $[x_0,b]$ é uma função derivável $y(x)$ que satisfaz a equação em
$[x_0,b]$ e a condição inicial.
\end{definicao}

\begin{teorema}[Existência e unicidade]\label{teo:picard}
Se $f$ é contínua em $D=\{(x,y):x_0\le x\le b,\ y\in\mathbb R\}$ e satisfaz a
\emph{condição de Lipschitz} em $y$,
\[
|f(x,y_1)-f(x,y_2)|\le L\,|y_1-y_2|\qquad\text{para todos }(x,y_1),(x,y_2)\in D,
\]
então o PVI tem uma única solução em $[x_0,b]$. A condição de Lipschitz vale, por exemplo,
quando $\partial f/\partial y$ é contínua e limitada em $D$, com
$L=\max|\partial f/\partial y|$.
\end{teorema}

Os métodos analíticos (variáveis separáveis, fator integrante, coeficientes constantes) só
resolvem uma pequena parte das EDOs que aparecem nas aplicações. Um método numérico não
produz uma fórmula para $y(x)$: ele produz \emph{aproximações} $y_i\approx y(x_i)$ numa malha
de pontos $x_i=x_0+ih$, $i=0,1,\dots,n$, em que $h$ é o \emph{passo}.

\paragraph{Interpretação geométrica: o campo de direções.} A equação $y'=f(x,y)$ diz qual é a
inclinação da solução que passa por cada ponto $(x,y)$ do plano. Desenhando pequenos segmentos
com essas inclinações, obtemos o \emph{campo de direções}; as soluções são as curvas que
seguem o campo (Figura~\ref{fig:campo}). Todos os métodos deste capítulo constroem a solução
aproximada andando passo a passo ``ao longo do campo''.

Usaremos ao longo do capítulo o PVI
\[
y'=y-x^2+1,\qquad y(0)=0{,}5,\qquad x\in[0,2],
\]
cuja solução exata, $y(x)=(x+1)^2-\tfrac12e^x$, nos permite medir os erros. Em particular,
$y(2)=9-\tfrac12e^2=5{,}3054719505\dots$

% ==================================================================
\section{Método de Euler}

\subsection{Dedução}

Pela fórmula de Taylor em torno de $x_i$, com $x_{i+1}=x_i+h$,
\[
y(x_{i+1})=y(x_i)+h\,y'(x_i)+\frac{h^2}{2}y''(\xi_i)
=y(x_i)+h\,f\bigl(x_i,y(x_i)\bigr)+\frac{h^2}{2}y''(\xi_i).
\]
Desprezando o termo em $h^2$ e trocando $y(x_i)$ pela aproximação $y_i$, obtemos o
\emph{método de Euler}:
\[
\boxed{\;y_{i+1}=y_i+h\,f(x_i,y_i),\qquad x_{i+1}=x_i+h.\;}
\]
Geometricamente, em cada passo seguimos a reta tangente à solução que passa por
$(x_i,y_i)$, durante um intervalo de comprimento $h$.

\begin{exemplo}[: o PVI do capítulo pelo método de Euler]
Com $h=0{,}2$ e $f(x,y)=y-x^2+1$:
\begin{align*}
y_1&=0{,}5+0{,}2\,(0{,}5-0+1)=0{,}5+0{,}2\cdot1{,}5=0{,}8,\\
y_2&=0{,}8+0{,}2\,(0{,}8-0{,}04+1)=0{,}8+0{,}2\cdot1{,}76=1{,}152,\\
y_3&=1{,}152+0{,}2\,(1{,}152-0{,}16+1)=1{,}152+0{,}2\cdot1{,}992=1{,}5504.
\end{align*}
Continuando até $x_{10}=2$:
\begin{center}
\renewcommand{\arraystretch}{1.25}
\begin{tabular}{c|c|c|c|c}
$i$ & $x_i$ & $y_i$ & $y(x_i)$ & $|y(x_i)-y_i|$\\ \hline
0 & 0 & $0{,}5$ & $0{,}5$ & 0\\
1 & $0{,}2$ & $0{,}8$ & $0{,}8292986$ & $0{,}0293$\\
2 & $0{,}4$ & $1{,}152$ & $1{,}2140877$ & $0{,}0621$\\
5 & $1{,}0$ & $2{,}4581760$ & $2{,}6408591$ & $0{,}1827$\\
10 & $2{,}0$ & $4{,}8657845$ & $5{,}3054720$ & $0{,}4397$
\end{tabular}
\end{center}
O erro cresce ao longo do intervalo: cada passo comete um pequeno erro e, além disso, parte
de um ponto que já estava errado.
\end{exemplo}

\begin{figure}[h]
\centering
\begin{tikzpicture}
\begin{axis}[width=10.5cm,height=7cm,axis lines=middle,xmin=-0.15,xmax=2.2,ymin=-0.2,ymax=6.2,
  xtick={0,0.5,1,1.5,2},xticklabels={$0$,$0{,}5$,$1$,$1{,}5$,$2$},grid=both,grid style={gray!15},
  xlabel={$x$},ylabel={$y$},unbounded coords=jump,
  legend style={at={(0.02,0.98)},anchor=north west,font=\small}]
\addplot[gray!60,thin,forget plot] coordinates {(-0.055,-0.055)(0.055,0.055)(nan,nan)(-0.045,0.433)(0.045,0.567)(nan,nan)(-0.037,0.926)(0.037,1.074)(nan,nan)(-0.031,1.422)(0.031,1.578)(nan,nan)(-0.027,1.919)(0.027,2.081)(nan,nan)(-0.024,2.418)(0.024,2.582)(nan,nan)(-0.021,2.916)(0.021,3.084)(nan,nan)(-0.019,3.416)(0.019,3.584)(nan,nan)(-0.017,3.915)(0.017,4.085)(nan,nan)(-0.016,4.415)(0.016,4.585)(nan,nan)(-0.014,4.914)(0.014,5.086)(nan,nan)(-0.013,5.414)(0.013,5.586)(nan,nan)(-0.012,5.914)(0.012,6.086)(nan,nan)(0.144,-0.053)(0.256,0.053)(nan,nan)(0.154,0.434)(0.246,0.566)(nan,nan)(0.162,0.926)(0.238,1.074)(nan,nan)(0.168,1.422)(0.232,1.578)(nan,nan)(0.173,1.919)(0.227,2.081)(nan,nan)(0.176,2.418)(0.224,2.582)(nan,nan)(0.179,2.917)(0.221,3.083)(nan,nan)(0.181,3.416)(0.219,3.584)(nan,nan)(0.183,3.915)(0.217,4.085)(nan,nan)(0.184,4.415)(0.216,4.585)(nan,nan)(0.186,4.914)(0.214,5.086)(nan,nan)(0.187,5.414)(0.213,5.586)(nan,nan)(0.188,5.914)(0.212,6.086)(nan,nan)(0.342,-0.049)(0.458,0.049)(nan,nan)(0.352,0.436)(0.448,0.564)(nan,nan)(0.361,0.928)(0.439,1.072)(nan,nan)(0.367,1.423)(0.433,1.577)(nan,nan)(0.372,1.920)(0.428,2.080)(nan,nan)(0.375,2.418)(0.425,2.582)(nan,nan)(0.378,2.917)(0.422,3.083)(nan,nan)(0.381,3.416)(0.419,3.584)(nan,nan)(0.382,3.915)(0.418,4.085)(nan,nan)(0.384,4.415)(0.416,4.585)(nan,nan)(0.385,4.914)(0.415,5.086)(nan,nan)(0.386,5.414)(0.414,5.586)(nan,nan)(0.387,5.914)(0.413,6.086)(nan,nan)(0.538,-0.040)(0.662,0.040)(nan,nan)(0.548,0.441)(0.652,0.559)(nan,nan)(0.558,0.930)(0.642,1.070)(nan,nan)(0.565,1.424)(0.635,1.576)(nan,nan)(0.570,1.921)(0.630,2.079)(nan,nan)(0.574,2.419)(0.626,2.581)(nan,nan)(0.577,2.917)(0.623,3.083)(nan,nan)(0.580,3.416)(0.620,3.584)(nan,nan)(0.582,3.916)(0.618,4.084)(nan,nan)(0.583,4.415)(0.617,4.585)(nan,nan)(0.585,4.915)(0.615,5.085)(nan,nan)(0.586,5.414)(0.614,5.586)(nan,nan)(0.587,5.914)(0.613,6.086)(nan,nan)(0.733,-0.024)(0.867,0.024)(nan,nan)(0.742,0.450)(0.858,0.550)(nan,nan)(0.753,0.936)(0.847,1.064)(nan,nan)(0.761,1.427)(0.839,1.573)(nan,nan)(0.767,1.923)(0.833,2.077)(nan,nan)(0.772,2.420)(0.828,2.580)(nan,nan)(0.776,2.918)(0.824,3.082)(nan,nan)(0.778,3.417)(0.822,3.583)(nan,nan)(0.781,3.916)(0.819,4.084)(nan,nan)(0.783,4.415)(0.817,4.585)(nan,nan)(0.784,4.915)(0.816,5.085)(nan,nan)(0.785,5.414)(0.815,5.586)(nan,nan)(0.787,5.914)(0.813,6.086)(nan,nan)(0.930,0.000)(1.070,0.000)(nan,nan)(0.935,0.468)(1.065,0.532)(nan,nan)(0.945,0.945)(1.055,1.055)(nan,nan)(0.955,1.433)(1.045,1.567)(nan,nan)(0.963,1.926)(1.037,2.074)(nan,nan)(0.969,2.422)(1.031,2.578)(nan,nan)(0.973,2.919)(1.027,3.081)(nan,nan)(0.976,3.418)(1.024,3.582)(nan,nan)(0.979,3.916)(1.021,4.084)(nan,nan)(0.981,4.416)(1.019,4.584)(nan,nan)(0.983,4.915)(1.017,5.085)(nan,nan)(0.984,5.415)(1.016,5.585)(nan,nan)(0.986,5.914)(1.014,6.086)(nan,nan)(1.134,0.029)(1.266,-0.029)(nan,nan)(1.130,0.496)(1.270,0.504)(nan,nan)(1.136,0.964)(1.264,1.036)(nan,nan)(1.147,1.443)(1.253,1.557)(nan,nan)(1.156,1.932)(1.244,2.068)(nan,nan)(1.164,2.425)(1.236,2.575)(nan,nan)(1.169,2.921)(1.231,3.079)(nan,nan)(1.174,3.419)(1.226,3.581)(nan,nan)(1.177,3.917)(1.223,4.083)(nan,nan)(1.179,4.416)(1.221,4.584)(nan,nan)(1.181,4.916)(1.219,5.084)(nan,nan)(1.183,5.415)(1.217,5.585)(nan,nan)(1.185,5.915)(1.215,6.085)(nan,nan)(1.344,0.053)(1.456,-0.053)(nan,nan)(1.334,0.530)(1.466,0.470)(nan,nan)(1.330,0.997)(1.470,1.003)(nan,nan)(1.336,1.465)(1.464,1.535)(nan,nan)(1.346,1.944)(1.454,2.056)(nan,nan)(1.356,2.432)(1.444,2.568)(nan,nan)(1.363,2.925)(1.437,3.075)(nan,nan)(1.369,3.421)(1.431,3.579)(nan,nan)(1.373,3.919)(1.427,4.081)(nan,nan)(1.377,4.417)(1.423,4.583)(nan,nan)(1.379,4.916)(1.421,5.084)(nan,nan)(1.381,5.416)(1.419,5.584)(nan,nan)(1.383,5.915)(1.417,6.085)(nan,nan)(1.556,0.068)(1.644,-0.068)(nan,nan)(1.547,0.557)(1.653,0.443)(nan,nan)(1.536,1.036)(1.664,0.964)(nan,nan)(1.530,1.504)(1.670,1.496)(nan,nan)(1.534,1.971)(1.666,2.029)(nan,nan)(1.544,2.447)(1.656,2.553)(nan,nan)(1.554,2.934)(1.646,3.066)(nan,nan)(1.562,3.426)(1.638,3.574)(nan,nan)(1.568,3.922)(1.632,4.078)(nan,nan)(1.573,4.419)(1.627,4.581)(nan,nan)(1.576,4.918)(1.624,5.082)(nan,nan)(1.579,5.417)(1.621,5.583)(nan,nan)(1.581,5.916)(1.619,6.084)(nan,nan)(1.766,0.076)(1.834,-0.076)(nan,nan)(1.759,0.571)(1.841,0.429)(nan,nan)(1.750,1.062)(1.850,0.938)(nan,nan)(1.740,1.545)(1.860,1.455)(nan,nan)(1.731,2.016)(1.869,1.984)(nan,nan)(1.731,2.482)(1.869,2.518)(nan,nan)(1.740,2.955)(1.860,3.045)(nan,nan)(1.751,3.438)(1.849,3.562)(nan,nan)(1.759,3.929)(1.841,4.071)(nan,nan)(1.766,4.423)(1.834,4.577)(nan,nan)(1.771,4.920)(1.829,5.080)(nan,nan)(1.775,5.418)(1.825,5.582)(nan,nan)(1.778,5.917)(1.822,6.083)(nan,nan)(1.973,0.081)(2.027,-0.081)(nan,nan)(1.969,0.578)(2.031,0.422)(nan,nan)(1.963,1.074)(2.037,0.926)(nan,nan)(1.955,1.567)(2.045,1.433)(nan,nan)(1.945,2.055)(2.055,1.945)(nan,nan)(1.935,2.532)(2.065,2.468)(nan,nan)(1.930,3.000)(2.070,3.000)(nan,nan)(1.935,3.468)(2.065,3.532)(nan,nan)(1.945,3.945)(2.055,4.055)(nan,nan)(1.955,4.433)(2.045,4.567)(nan,nan)(1.963,4.926)(2.037,5.074)(nan,nan)(1.969,5.422)(2.031,5.578)(nan,nan)(1.973,5.919)(2.027,6.081)(nan,nan)};
\addplot[domain=0:2,samples=80,very thick,verde] {(x+1)^2-0.5*exp(x)};
\addplot[thick,vermelho,mark=*,mark size=1.6pt] coordinates {(0.00,0.50000)(0.20,0.80000)(0.40,1.15200)(0.60,1.55040)(0.80,1.98848)(1.00,2.45818)(1.20,2.94981)(1.40,3.45177)(1.60,3.95013)(1.80,4.42815)(2.00,4.86578)};
\legend{solução exata,Euler com $h=0{,}2$}
\end{axis}
\end{tikzpicture}
\caption{Campo de direções de $y'=y-x^2+1$, a solução com $y(0)=0{,}5$ e a poligonal de
Euler com $h=0{,}2$.}
\label{fig:campo}
\end{figure}

\subsection{Erro local e erro global}

Dois tipos de erro aparecem em qualquer método de passo a passo:
\begin{itemize}[nosep]
  \item o \emph{erro local de truncamento} é o erro cometido num único passo, supondo que se
        parte do valor exato $y(x_i)$. No método de Euler, pela fórmula de Taylor acima, ele é
        $\frac{h^2}{2}y''(\xi_i)$, isto é, $O(h^2)$;
  \item o \emph{erro global} $e_i=y(x_i)-y_i$ acumula os erros locais dos $i$ passos e a
        propagação dos erros anteriores.
\end{itemize}
Como o número de passos até um ponto fixo $\bar x$ é $(\bar x-x_0)/h$, somar erros locais
$O(h^2)$ dá um erro global $O(h)$. O teorema a seguir torna isso preciso.

\begin{teorema}[Erro global do método de Euler]\label{teo:erroeuler}
Suponha $f$ nas condições do Teorema~\ref{teo:picard}, com constante de Lipschitz $L$, e
$|y''(x)|\le M$ em $[x_0,b]$. Então as aproximações de Euler satisfazem
\[
|y(x_i)-y_i|\le\frac{hM}{2L}\Bigl(e^{L(x_i-x_0)}-1\Bigr),\qquad i=0,1,\dots,n.
\]
\end{teorema}

\begin{proof}[Ideia da demonstração]
Subtraindo $y_{i+1}=y_i+hf(x_i,y_i)$ de
$y(x_{i+1})=y(x_i)+hf(x_i,y(x_i))+\frac{h^2}{2}y''(\xi_i)$ e usando a condição de
Lipschitz, $|e_{i+1}|\le(1+hL)|e_i|+\frac{h^2M}{2}$. Por indução,
$|e_i|\le\frac{hM}{2L}\bigl((1+hL)^i-1\bigr)$, e $(1+hL)^i\le e^{ihL}=e^{L(x_i-x_0)}$.
\end{proof}

A cota é proporcional a $h$: o método de Euler é de \emph{primeira ordem}. Dividir o passo por
2 divide o erro (aproximadamente) por 2; para ganhar uma casa decimal é preciso usar 10 vezes
mais passos. O fator exponencial mostra também que o erro pode crescer muito em intervalos
longos.

\begin{atividade}{método de Euler}{\arqEuler}
Abra o aplicativo. À esquerda ficam o campo de direções, a poligonal de Euler (vermelha) e a
solução de referência (verde); à direita, a fórmula, o erro em $x_n$ e a tabela com $x_i$,
$y_i$, $f(x_i,y_i)$, o valor de referência e o erro em cada ponto.
\begin{enumerate}[nosep]
  \item Confira os números do exemplo com $h=0{,}2$ e $n=10$.
  \item Use $h=0{,}1$ e $n=20$ (mesmo ponto final $x_n=2$). O erro em $x_n$ caiu
        aproximadamente pela metade?
  \item Marque \emph{Soluções exatas que passam por cada ponto}. Cada passo de Euler leva a
        aproximação para uma curva solução diferente. Use essa figura para explicar a
        diferença entre erro local e erro global.
  \item Digite $f(x,y)=-20y$, com $y_0=1$, $h=0{,}12$ e $n=10$. O que acontece? Repita com
        $h=0{,}04$. (Voltaremos a isso na Seção~\ref{sec:estab}.)
\end{enumerate}
\end{atividade}

% ==================================================================
\section{Métodos de Runge-Kutta}

\subsection{A ideia}

Para aumentar a ordem poderíamos usar mais termos da série de Taylor,
$y_{i+1}=y_i+hy'_i+\frac{h^2}{2}y''_i+\cdots$ (\emph{métodos de Taylor}), mas isso exige
calcular derivadas de $f$: no exemplo, $y''=\frac{d}{dx}(y-x^2+1)=y'-2x=y-x^2+1-2x$. Para
uma $f$ complicada, essas contas são trabalhosas e precisam ser refeitas para cada problema.

Os \emph{métodos de Runge-Kutta} obtêm a mesma ordem usando apenas valores de $f$: em cada
passo, $f$ é avaliada em alguns pontos de teste e as inclinações obtidas são combinadas com
pesos escolhidos para que a fórmula concorde com a série de Taylor até a ordem desejada.
Euler é o método de Runge-Kutta de ordem 1, com uma única inclinação.

\subsection{Métodos de segunda ordem}

Procuramos uma fórmula
\[
y_{i+1}=y_i+h\,(a_1k_1+a_2k_2),\qquad k_1=f(x_i,y_i),\quad k_2=f(x_i+p\,h,\ y_i+p\,h\,k_1).
\]
Expandindo $k_2$ em série de Taylor em duas variáveis e comparando com
$y(x_i+h)=y+hf+\frac{h^2}{2}(f_x+f_yf)+O(h^3)$, obtemos as condições
\[
a_1+a_2=1,\qquad a_2\,p=\tfrac12.
\]
Há infinitas escolhas; as duas mais usadas são:
\begin{itemize}[nosep]
  \item \textbf{método de Heun} (ou \emph{Euler melhorado}), com $a_1=a_2=\frac12$ e $p=1$:
        \[
        k_1=f(x_i,y_i),\quad k_2=f(x_i+h,\ y_i+hk_1),\quad y_{i+1}=y_i+\frac h2(k_1+k_2);
        \]
        primeiro damos um passo de Euler (\emph{previsão}) e depois corrigimos com a média das
        inclinações no início e no fim (\emph{correção}). Se $f$ não depende de $y$, isso é
        exatamente a regra dos trapézios do Capítulo~\ref{cap:integracao};
  \item \textbf{método do ponto médio}, com $a_1=0$, $a_2=1$ e $p=\frac12$:
        \[
        k_1=f(x_i,y_i),\quad k_2=f\bigl(x_i+\tfrac h2,\ y_i+\tfrac h2k_1\bigr),\quad y_{i+1}=y_i+h\,k_2;
        \]
        usamos a inclinação estimada no meio do passo.
\end{itemize}
Ambos têm erro local $O(h^3)$ e erro global $O(h^2)$, com duas avaliações de $f$ por passo.

\begin{exemplo}[: primeiro passo de Heun e do ponto médio]
Com $h=0{,}2$, $k_1=f(0;\,0{,}5)=1{,}5$ nos dois métodos.

\textbf{Heun:} o passo de Euler leva a $(0{,}2;\ 0{,}8)$ e
$k_2=f(0{,}2;\,0{,}8)=0{,}8-0{,}04+1=1{,}76$; logo
$y_1=0{,}5+0{,}1\,(1{,}5+1{,}76)=0{,}826$.

\textbf{Ponto médio:} o meio passo leva a $(0{,}1;\ 0{,}65)$ e
$k_2=f(0{,}1;\,0{,}65)=0{,}65-0{,}01+1=1{,}64$; logo $y_1=0{,}5+0{,}2\cdot1{,}64=0{,}828$.

O valor exato é $y(0{,}2)=0{,}8292986$. Prosseguindo até $x=2$, Heun dá
$y_{10}=5{,}2330546$ (erro $0{,}0724$) e o ponto médio dá $y_{10}=5{,}2903695$ (erro
$0{,}0151$), contra $0{,}4397$ de Euler.
\end{exemplo}

\begin{atividade}{Runge-Kutta de segunda ordem}{\arqRKdois}
O aplicativo mostra, para o passo escolhido no controle $kk$, a inclinação $k_1$ (laranja), o
ponto de teste ($\times$) e a inclinação $k_2$ (roxa). A caixa \emph{Usar o ponto médio}
troca Heun pelo método do ponto médio.
\begin{enumerate}[nosep]
  \item Confira o primeiro passo do exemplo nos dois métodos ($kk=0$).
  \item Marque \emph{Comparar com Euler} e compare as duas poligonais.
  \item Com Heun, use $h=0{,}2$, $n=10$ e depois $h=0{,}1$, $n=20$. O erro em $x_n$ fica
        cerca de 4 vezes menor? Por quê?
  \item Digite $f(x,y)=\cos(x)$, com $x_0=0$ e $y_0=0$. Que regra de integração numérica o
        método de Heun está aplicando? E o método do ponto médio?
\end{enumerate}
\end{atividade}

\subsection{O método de Runge-Kutta de quarta ordem}

O método mais usado na prática combina quatro inclinações por passo:
\[
\boxed{\;
\begin{aligned}
k_1&=f(x_i,\,y_i), & k_2&=f\bigl(x_i+\tfrac h2,\ y_i+\tfrac h2k_1\bigr),\\
k_3&=f\bigl(x_i+\tfrac h2,\ y_i+\tfrac h2k_2\bigr), & k_4&=f(x_i+h,\ y_i+hk_3),\\
y_{i+1}&=y_i+\frac h6\,(k_1+2k_2+2k_3+k_4). & &
\end{aligned}\;}
\]
A inclinação usada é uma média ponderada: uma no início do passo, duas no meio (pesos maiores)
e uma no fim. Quando $f$ só depende de $x$, $k_2=k_3=f(x_i+\frac h2)$ e a fórmula fica
$\frac h6\bigl[f(x_i)+4f(x_i+\frac h2)+f(x_i+h)\bigr]$: é a regra $1/3$ de Simpson no intervalo
$[x_i,x_{i+1}]$. O RK4 tem erro local $O(h^5)$ e erro global $O(h^4)$.

\begin{exemplo}[: o PVI do capítulo pelo RK4]
Primeiro passo, com $h=0{,}2$:
\begin{align*}
k_1&=f(0;\ 0{,}5)=1{,}5,\\
k_2&=f(0{,}1;\ 0{,}5+0{,}1\cdot1{,}5)=f(0{,}1;\ 0{,}65)=1{,}64,\\
k_3&=f(0{,}1;\ 0{,}5+0{,}1\cdot1{,}64)=f(0{,}1;\ 0{,}664)=1{,}654,\\
k_4&=f(0{,}2;\ 0{,}5+0{,}2\cdot1{,}654)=f(0{,}2;\ 0{,}8308)=1{,}7908,\\
y_1&=0{,}5+\frac{0{,}2}{6}\,(1{,}5+2\cdot1{,}64+2\cdot1{,}654+1{,}7908)=0{,}8292933,
\end{align*}
com erro $5{,}3\times10^{-6}$. Ao final, em $x=2$, $y_{10}=5{,}3053630$, com erro
$1{,}09\times10^{-4}$, cerca de 4000 vezes menor que o de Euler com o mesmo passo.
\end{exemplo}

\begin{atividade}{Runge-Kutta de quarta ordem}{\arqRKquatro}
Para o passo escolhido em $kk$, o gráfico mostra as quatro inclinações como pequenos segmentos
desenhados nos pontos em que foram calculadas: $k_1$ no ponto atual, $k_2$ e $k_3$ no meio do
passo e $k_4$ no fim. A tabela traz $k_1,\dots,k_4$ e o erro de cada $y_i$.
\begin{enumerate}[nosep]
  \item Confira os números do primeiro passo ($kk=0$) e o valor final.
  \item Compare os erros com $h=0{,}2$ ($n=10$) e $h=0{,}1$ ($n=20$). Qual é a razão?
  \item Use $f(x,y)=\cos(x)$, $x_0=0$, $y_0=0$, $h=0{,}5$ e $n=4$. Compare $y_4$ com
        $\operatorname{sen}2=\int_0^2\cos x\,dx$ e com a regra $1/3$ de Simpson com 8
        subintervalos do Capítulo~\ref{cap:integracao}.
\end{enumerate}
\end{atividade}

% ==================================================================
\section{Ordem, custo e estabilidade}\label{sec:estab}

\subsection{Ordem e custo}

Dizemos que um método tem \emph{ordem $p$} se o erro global em um ponto fixo é $O(h^p)$. Na
prática, isso significa que, ao dividir $h$ por 2, o erro fica aproximadamente $2^p$ vezes
menor. A tabela mostra os erros em $x=2$ para o problema do capítulo.

\begin{center}
\renewcommand{\arraystretch}{1.3}
\begin{tabular}{c|c|cc|cc|cc}
$n$ & $h$ & erro Euler & razão & erro Heun & razão & erro RK4 & razão\\ \hline
5 & $0{,}4$ & $7{,}46\times10^{-1}$ & & $2{,}63\times10^{-1}$ & & $1{,}64\times10^{-3}$ & \\
10 & $0{,}2$ & $4{,}40\times10^{-1}$ & $1{,}70$ & $7{,}24\times10^{-2}$ & $3{,}63$ & $1{,}09\times10^{-4}$ & $15{,}1$\\
20 & $0{,}1$ & $2{,}42\times10^{-1}$ & $1{,}82$ & $1{,}89\times10^{-2}$ & $3{,}83$ & $6{,}99\times10^{-6}$ & $15{,}6$\\
40 & $0{,}05$ & $1{,}27\times10^{-1}$ & $1{,}90$ & $4{,}82\times10^{-3}$ & $3{,}92$ & $4{,}42\times10^{-7}$ & $15{,}8$\\
80 & $0{,}025$ & $6{,}55\times10^{-2}$ & $1{,}95$ & $1{,}22\times10^{-3}$ & $3{,}96$ & $2{,}78\times10^{-8}$ & $15{,}9$
\end{tabular}
\end{center}

As razões se aproximam de $2$, $4$ e $16$, confirmando as ordens 1, 2 e 4. O custo de um
passo é medido pelo número de avaliações de $f$: 1 para Euler, 2 para Heun e 4 para RK4.
Mesmo assim o RK4 compensa: com $n=10$ (40 avaliações) ele é muito mais preciso que Euler com
$n=80$ (80 avaliações).

\begin{figure}[h]
\centering
\begin{tikzpicture}
\begin{loglogaxis}[width=10cm,height=6.5cm,xlabel={passo $h$},ylabel={erro em $x=2$},
  grid=both,grid style={gray!20},x dir=reverse,xtick={0.025,0.05,0.1,0.2,0.4},xticklabels={$0{,}025$,$0{,}05$,$0{,}1$,$0{,}2$,$0{,}4$},
  legend style={at={(0.02,0.02)},anchor=south west,font=\small}]
\addplot[thick,cL1,mark=*] coordinates {(0.4,0.746)(0.2,0.440)(0.1,0.242)(0.05,0.127)(0.025,0.0655)};
\addplot[thick,cL3,mark=square*] coordinates {(0.4,0.263)(0.2,0.0724)(0.1,0.0189)(0.05,0.00482)(0.025,0.00122)};
\addplot[thick,vermelho,mark=triangle*] coordinates {(0.4,1.64e-3)(0.2,1.09e-4)(0.1,6.99e-6)(0.05,4.42e-7)(0.025,2.78e-8)};
\legend{Euler (ordem 1),Heun (ordem 2),RK4 (ordem 4)}
\end{loglogaxis}
\end{tikzpicture}
\caption{Erro em $x=2$ em função do passo, em escala logarítmica. As inclinações das retas
são as ordens dos métodos.}
\label{fig:ordem}
\end{figure}

\begin{atividade}{ordem dos métodos}{\arqOrdem}
O aplicativo resolve o mesmo PVI por Euler, Heun e RK4 com $n_0$, $2n_0$, $4n_0$ e $8n_0$
passos e mostra os erros e as razões entre erros consecutivos.
\begin{enumerate}[nosep]
  \item Com $n_0=5$, confira a tabela acima.
  \item Digite outro PVI, por exemplo $f(x,y)=-2xy$ com $y_0=1$ e $x_f=1$ (solução $e^{-x^2}$),
        e verifique que as razões continuam próximas de 2, 4 e 16.
  \item Compare o erro do RK4 com $n_0$ passos e o de Euler com $8n_0$ passos (o dobro do custo
        do RK4). Qual método compensa?
\end{enumerate}
\end{atividade}

\subsection{Estabilidade}

Um método convergente pode, mesmo assim, dar resultados absurdos se o passo for grande demais.
O exemplo-modelo é
\[
y'=\lambda y,\qquad y(0)=1,\qquad \lambda<0,
\]
cuja solução $y=e^{\lambda x}$ decai para zero. Pelo método de Euler,
$y_{i+1}=y_i+h\lambda y_i=(1+h\lambda)\,y_i$, logo $y_i=(1+h\lambda)^i$. A aproximação só
decai se $|1+h\lambda|<1$, isto é,
\[
h<\frac{2}{|\lambda|}.
\]
Com $\lambda=-20$ e $h=0{,}12$, o fator é $1+h\lambda=-1{,}4$ e as aproximações oscilam com
amplitude crescente: $-1{,}4$; $1{,}96$; $-2{,}744$; $3{,}84$; $\dots$ Com $h=0{,}04$, o fator é $0{,}2$
e elas decaem corretamente. Esse fenômeno se chama \emph{instabilidade numérica}: não é um
erro de arredondamento, é a própria fórmula que amplifica as perturbações.

Os métodos de Runge-Kutta também têm um limite de estabilidade (para o RK4, $h|\lambda|\lesssim2{,}78$
quando $\lambda$ é real e negativo). Problemas em que convivem escalas de tempo muito diferentes
(variações rápidas que somem logo e variações lentas que interessam) são chamados
\emph{rígidos} (\emph{stiff}); para eles usam-se métodos implícitos, que não têm essa
restrição. Esse assunto vai além deste livro.

% ==================================================================
\section{Sistemas e equações de ordem superior}

\subsection{Sistemas de EDOs}

Muitos modelos envolvem várias grandezas que interagem. Um sistema de duas EDOs de primeira
ordem
\[
\begin{cases}u'=F(x,u,v),& u(x_0)=u_0,\\ v'=G(x,u,v),& v(x_0)=v_0,\end{cases}
\]
pode ser escrito na forma vetorial $\mathbf Y'=\mathbf F(x,\mathbf Y)$, com $\mathbf Y=(u,v)$ e
$\mathbf F=(F,G)$. Todos os métodos anteriores valem \emph{sem mudança}, desde que as contas
sejam feitas com vetores. Para o RK4:
\begin{align*}
\mathbf K_1&=\mathbf F(x_i,\mathbf Y_i), &
\mathbf K_2&=\mathbf F\bigl(x_i+\tfrac h2,\ \mathbf Y_i+\tfrac h2\mathbf K_1\bigr),\\
\mathbf K_3&=\mathbf F\bigl(x_i+\tfrac h2,\ \mathbf Y_i+\tfrac h2\mathbf K_2\bigr), &
\mathbf K_4&=\mathbf F(x_i+h,\ \mathbf Y_i+h\mathbf K_3),
\end{align*}
\[
\mathbf Y_{i+1}=\mathbf Y_i+\frac h6\,(\mathbf K_1+2\mathbf K_2+2\mathbf K_3+\mathbf K_4).
\]
Cuidado: $\mathbf K_2$ só pode ser calculado depois de \emph{todas} as componentes de
$\mathbf K_1$, e assim por diante.

\subsection{Equações de ordem superior}

Uma EDO de segunda ordem $y''=g(x,y,y')$, com $y(x_0)=y_0$ e $y'(x_0)=z_0$, vira um sistema de
primeira ordem com a troca $u=y$, $v=y'$:
\[
u'=v,\qquad v'=g(x,u,v),\qquad u(x_0)=y_0,\quad v(x_0)=z_0.
\]
Por exemplo, o pêndulo simples $\theta''=-\frac gL\operatorname{sen}\theta$ vira
$u'=v$, $v'=-\frac gL\operatorname{sen}u$. O mesmo vale para ordem $m$: introduzimos $y,y',\dots,y^{(m-1)}$
como incógnitas.

\begin{exemplo}[: o modelo predador-presa de Lotka-Volterra]
Seja $u(t)$ uma população de presas e $v(t)$ uma população de predadores (em unidades
convenientes), com
\[
u'=1{,}1\,u-0{,}4\,uv,\qquad v'=0{,}1\,uv-0{,}4\,v,\qquad u(0)=10,\quad v(0)=5.
\]
Sem predadores, as presas cresceriam exponencialmente; sem presas, os predadores morreriam.
Pode-se mostrar que as soluções são \emph{periódicas}: no plano $(u,v)$ (o \emph{plano de
fase}) elas percorrem curvas fechadas. Com $h=0{,}1$, o RK4 reproduz a órbita fechada; o
método de Euler, com o mesmo passo, gera uma espiral que se afasta dela
(Figura~\ref{fig:lv}): em $t=20$, Euler dá $(u,v)\approx(3{,}32;\ 0{,}50)$, enquanto o RK4 dá
$(10{,}21;\ 1{,}33)$. Mesmo com $h=0{,}01$, Euler ainda erra visivelmente.
\end{exemplo}

\begin{figure}[h]
\centering
\begin{tikzpicture}
\begin{axis}[width=10cm,height=6.5cm,xlabel={$u$ (presas)},ylabel={$v$ (predadores)},
  grid=both,grid style={gray!20},xmin=0,xmax=15,ymin=0,ymax=7.5,
  legend style={at={(0.98,0.98)},anchor=north east,font=\small}]
\addplot[thick,vermelho] coordinates {(10.000,5.000)(8.168,5.535)(6.432,5.911)(4.949,6.111)(3.772,6.152)(2.881,6.067)(2.224,5.892)(1.745,5.658)(1.398,5.389)(1.145,5.102)(0.960,4.809)(0.823,4.519)(0.723,4.237)(0.649,3.965)(0.595,3.706)(0.557,3.460)(0.531,3.229)(0.515,3.012)(0.509,2.809)(0.510,2.620)(0.519,2.443)(0.536,2.279)(0.560,2.127)(0.592,1.986)(0.632,1.856)(0.683,1.736)(0.744,1.626)(0.817,1.524)(0.905,1.431)(1.009,1.347)(1.132,1.270)(1.278,1.201)(1.450,1.139)(1.653,1.085)(1.892,1.038)(2.174,0.997)(2.505,0.965)(2.892,0.940)(3.345,0.923)(3.873,0.916)(4.485,0.919)(5.189,0.935)(5.994,0.965)(6.902,1.013)(7.910,1.084)(9.003,1.185)(10.150,1.325)(11.294,1.515)(12.345,1.772)(13.176,2.113)(13.634,2.552)(13.563,3.095)(12.871,3.725)(11.590,4.396)(9.900,5.033)(8.067,5.561)(6.342,5.927)(4.875,6.117)(3.716,6.151)(2.839,6.060)(2.193,5.881)(1.723,5.644)(1.382,5.373)(1.133,5.086)(0.951,4.793)(0.817,4.503)(0.718,4.221)(0.645,3.950)(0.592,3.692)(0.555,3.447)(0.530,3.217)(0.515,3.000)(0.509,2.798)(0.511,2.610)(0.520,2.434)(0.537,2.271)(0.561,2.119)(0.594,1.979)(0.635,1.849)(0.686,1.730)(0.747,1.620)(0.821,1.519)(0.910,1.427)(1.015,1.342)(1.140,1.266)(1.287,1.197)(1.461,1.136)(1.666,1.082)(1.907,1.035)(2.191,0.995)(2.524,0.963)(2.915,0.939)(3.372,0.923)(3.905,0.916)(4.521,0.920)(5.231,0.936)(6.042,0.967)(6.955,1.016)(7.968,1.089)(9.066,1.192)(10.214,1.334)};
\addplot[thick,gray,dashed] coordinates {(10.000,5.000)(8.172,5.570)(6.365,5.991)(4.790,6.227)(3.542,6.284)(2.615,6.196)(1.951,6.003)(1.485,5.745)(1.158,5.449)(0.928,5.137)(0.764,4.822)(0.647,4.512)(0.563,4.213)(0.502,3.927)(0.459,3.656)(0.428,3.401)(0.408,3.162)(0.397,2.939)(0.392,2.731)(0.395,2.537)(0.403,2.358)(0.417,2.191)(0.437,2.037)(0.464,1.895)(0.497,1.764)(0.538,1.642)(0.588,1.531)(0.649,1.429)(0.720,1.335)(0.805,1.250)(0.906,1.172)(1.026,1.101)(1.167,1.037)(1.334,0.980)(1.531,0.929)(1.764,0.885)(2.038,0.847)(2.361,0.815)(2.742,0.790)(3.190,0.772)(3.716,0.762)(4.331,0.760)(5.048,0.767)(5.878,0.787)(6.834,0.820)(7.923,0.872)(9.147,0.947)(10.495,1.053)(11.937,1.203)(13.412,1.411)(14.812,1.700)(15.969,2.099)(16.647,2.641)(16.569,3.354)(15.506,4.235)(13.435,5.222)(10.671,6.178)(7.794,6.933)(5.335,7.375)(3.527,7.500)(2.318,7.379)(1.549,7.101)(1.065,6.739)(0.760,6.338)(0.563,5.927)(0.433,5.523)(0.345,5.133)(0.285,4.763)(0.243,4.415)(0.213,4.089)(0.192,3.784)(0.177,3.501)(0.168,3.238)(0.162,2.995)(0.160,2.769)(0.160,2.561)(0.163,2.368)(0.169,2.190)(0.177,2.025)(0.188,1.873)(0.202,1.734)(0.219,1.604)(0.240,1.486)(0.266,1.376)(0.297,1.276)(0.333,1.183)(0.377,1.098)(0.429,1.020)(0.491,0.949)(0.565,0.884)(0.653,0.825)(0.759,0.771)(0.884,0.722)(1.035,0.678)(1.214,0.639)(1.429,0.605)(1.685,0.575)(1.992,0.549)(2.359,0.528)(2.798,0.512)(3.322,0.501)};
\addplot[only marks,mark=*,azul] coordinates {(10,5)};
\legend{RK4 ($h=0{,}1$),Euler ($h=0{,}1$),{$(u_0,v_0)$}}
\end{axis}
\end{tikzpicture}
\caption{Plano de fase do modelo de Lotka-Volterra para $0\le t\le20$.}
\label{fig:lv}
\end{figure}

\begin{atividade}{sistemas de EDOs com RK4}{\arqSistemas}
O aplicativo resolve um sistema autônomo $u'=F(u,v)$, $v'=G(u,v)$ pelo RK4. As funções $F$ e
$G$ podem ser digitadas; a caixa \emph{Gráficos u(t) e v(t)} troca o plano de fase pelos
gráficos em função do tempo.
\begin{enumerate}[nosep]
  \item Confira os valores do exemplo em $t=20$ ($h=0{,}1$, $n=200$), com e sem \emph{Comparar
        com Euler}.
  \item Nos gráficos $u(t)$ e $v(t)$, estime o período das oscilações. Por que o pico dos
        predadores vem depois do pico das presas?
  \item Digite $F=v$ e $G=-u$ (oscilador harmônico $y''=-y$), com $u_0=1$, $v_0=0$. A órbita
        exata é a circunferência $u^2+v^2=1$. Compare RK4 e Euler com $h=0{,}1$ e $n=200$.
  \item Pêndulo: $F=v$ e $G=-9{,}8\operatorname{sen}(u)$ (digite \texttt{-9.8 sin(u)}), com
        $u_0=1$ e $v_0=0$. Compare com a aproximação linear $G=-9{,}8u$.
\end{enumerate}
\end{atividade}

% ==================================================================
\section{Problemas de valor de contorno}

Num PVI, todas as condições são dadas no mesmo ponto. Em muitos problemas, porém, as condições
aparecem nos dois extremos de um intervalo: a deflexão de uma viga apoiada nas pontas, a
temperatura numa barra com as extremidades mantidas a temperaturas fixas, a forma de um cabo
suspenso.

\begin{definicao}
Um \emph{problema de valor de contorno} (PVC) de segunda ordem é
\[
y''=g(x,y,y'),\qquad a\le x\le b,\qquad y(a)=\alpha,\quad y(b)=\beta.
\]
\end{definicao}

Diferentemente do PVI, um PVC pode não ter solução ou ter infinitas: $y''=-y$,
$y(0)=0$, $y(\pi)=1$ não tem solução (todas as soluções com $y(0)=0$ são $c\operatorname{sen}x$
e valem 0 em $\pi$), enquanto $y''=-y$, $y(0)=0$, $y(\pi)=0$ tem infinitas. Para o caso
linear há um resultado simples:

\begin{teorema}\label{teo:pvclinear}
Se $p$, $q$ e $r$ são contínuas em $[a,b]$ e $q(x)>0$ em $[a,b]$, o PVC linear
\[
y''=p(x)\,y'+q(x)\,y+r(x),\qquad y(a)=\alpha,\quad y(b)=\beta,
\]
tem uma única solução.
\end{teorema}

Estudaremos dois métodos: o \emph{método do tiro}, que transforma o PVC em PVIs (e reaproveita
o RK4), e o \emph{método das diferenças finitas}, que transforma o PVC num sistema linear (e
reaproveita o Capítulo~\ref{cap:sistemas}).

% ==================================================================
\section{Método do tiro}

\subsection{A ideia}

Se conhecêssemos $y'(a)$, o PVC seria um PVI. O método do tiro ``chuta'' esse valor:
para cada $s$, resolvemos
\[
y''=g(x,y,y'),\qquad y(a)=\alpha,\qquad y'(a)=s,
\]
por RK4 (como sistema de primeira ordem) e olhamos onde a solução, que chamamos $y(x;s)$,
chega em $x=b$. Como um canhão que corrige a inclinação depois de cada tiro, ajustamos $s$
até acertar o alvo $\beta$. Em termos do Capítulo~\ref{cap:zeros}, procuramos um zero da função
\[
\phi(s)=y(b;s)-\beta.
\]
Cada avaliação de $\phi$ exige resolver um PVI inteiro, e não temos uma fórmula para
$\phi'(s)$; por isso o \textbf{método da secante} é a escolha natural:
\[
s_{k+1}=s_k-\phi(s_k)\,\frac{s_k-s_{k-1}}{\phi(s_k)-\phi(s_{k-1})}.
\]

\begin{observacao}
Se a EDO é \emph{linear}, $y(b;s)$ é uma função afim de $s$ (combinação de duas soluções, por
superposição). Então $\phi$ é uma reta e a secante encontra o zero exatamente no primeiro
passo: bastam dois tiros e uma combinação.
\end{observacao}

\begin{exemplo}[: um PVC não linear]
Considere
\[
y''=\frac{1}{8}\bigl(32+2x^3-y\,y'\bigr),\qquad 1\le x\le3,\qquad y(1)=17,\quad y(3)=\tfrac{43}{3},
\]
cuja solução exata é $y=x^2+\frac{16}{x}$ (assim, $y'(1)=-14$). Usando RK4 com $h=0{,}1$
(20 passos) em cada tiro, os chutes $s_0=-1{,}4$ e $s_1=s_0+1=-0{,}4$ e a secante:
\begin{center}
\renewcommand{\arraystretch}{1.25}
\begin{tabular}{c|r|r|r}
$k$ & $s_k$ & $y(3;s_k)$ & $\phi(s_k)$\\ \hline
0 & $-1{,}4$ & $20{,}4518393$ & $6{,}12$\\
1 & $-0{,}4$ & $20{,}8582730$ & $6{,}52$\\
2 & $-16{,}4541312$ & $12{,}8171044$ & $-1{,}52$\\
3 & $-13{,}4269919$ & $14{,}6667588$ & $3{,}3\times10^{-1}$\\
4 & $-13{,}9726751$ & $14{,}3495052$ & $1{,}6\times10^{-2}$\\
5 & $-14{,}0004910$ & $14{,}3331574$ & $-1{,}8\times10^{-4}$\\
6 & $-14{,}0001918$ & $14{,}3333334$ & $9{,}2\times10^{-8}$\\
7 & $-14{,}0001919$ & $14{,}3333333$ & $5\times10^{-13}$
\end{tabular}
\end{center}
O valor final de $s$ difere de $-14$ em $2\times10^{-4}$: o que resta é o erro do RK4 com
$h=0{,}1$, e não o da secante. Com esse $s$, o erro máximo da solução em $[1,3]$ é da ordem
de $10^{-4}$.
\end{exemplo}

\begin{figure}[h]
\centering
\begin{tikzpicture}
\begin{axis}[width=10cm,height=6.5cm,xlabel={$x$},ylabel={$y$},grid=both,grid style={gray!20},
  xmin=0.9,xmax=3.1,ymin=9.5,ymax=23,legend style={at={(0.02,0.98)},anchor=north west,font=\small},legend columns=2]
\addplot[thick,gray] coordinates {(1.00,17.0000)(1.10,16.8938)(1.20,16.8469)(1.30,16.8489)(1.40,16.8916)(1.50,16.9684)(1.60,17.0743)(1.70,17.2050)(1.80,17.3576)(1.90,17.5295)(2.00,17.7189)(2.10,17.9246)(2.20,18.1457)(2.30,18.3815)(2.40,18.6319)(2.50,18.8969)(2.60,19.1766)(2.70,19.4715)(2.80,19.7819)(2.90,20.1085)(3.00,20.4518)};
\addplot[thick,cL1] coordinates {(1.00,17.0000)(1.10,16.9839)(1.20,17.0100)(1.30,17.0710)(1.40,17.1613)(1.50,17.2763)(1.60,17.4126)(1.70,17.5675)(1.80,17.7387)(1.90,17.9250)(2.00,18.1251)(2.10,18.3383)(2.20,18.5644)(2.30,18.8032)(2.40,19.0548)(2.50,19.3195)(2.60,19.5977)(2.70,19.8900)(2.80,20.1971)(2.90,20.5196)(3.00,20.8583)};
\addplot[thick,cL3] coordinates {(1.00,17.0000)(1.10,15.5332)(1.20,14.3654)(1.30,13.4307)(1.40,12.6826)(1.50,12.0878)(1.60,11.6216)(1.70,11.2653)(1.80,11.0047)(1.90,10.8287)(2.00,10.7285)(2.10,10.6971)(2.20,10.7288)(2.30,10.8189)(2.40,10.9635)(2.50,11.1593)(2.60,11.4035)(2.70,11.6935)(2.80,12.0271)(2.90,12.4022)(3.00,12.8171)};
\addplot[thick,cL4] coordinates {(1.00,17.0000)(1.10,15.8074)(1.20,14.8685)(1.30,14.1297)(1.40,13.5524)(1.50,13.1084)(1.60,12.7763)(1.70,12.5397)(1.80,12.3860)(1.90,12.3049)(2.00,12.2885)(2.10,12.3300)(2.20,12.4242)(2.30,12.5666)(2.40,12.7535)(2.50,12.9819)(2.60,13.2490)(2.70,13.5528)(2.80,13.8914)(2.90,14.2632)(3.00,14.6668)};
\addplot[very thick,vermelho] coordinates {(1.00,17.0000)(1.10,15.7555)(1.20,14.7734)(1.30,13.9978)(1.40,13.3886)(1.50,12.9167)(1.60,12.5601)(1.70,12.3018)(1.80,12.1289)(1.90,12.0311)(2.00,12.0000)(2.10,12.0291)(2.20,12.1127)(2.30,12.2465)(2.40,12.4267)(2.50,12.6500)(2.60,12.9139)(2.70,13.2159)(2.80,13.5543)(2.90,13.9272)(3.00,14.3333)};
\addplot[only marks,mark=*,mark size=2.5pt] coordinates {(1,17)(3,14.3333)};
\legend{$s_0$,$s_1$,$s_2$,$s_3$,$s_7$}
\end{axis}
\end{tikzpicture}
\caption{Método do tiro: as soluções dos PVIs com $y'(1)=s_k$. O alvo é o ponto $(3,\frac{43}{3})$.}
\label{fig:tiro}
\end{figure}

\begin{atividade}{método do tiro}{\arqTiro}
O aplicativo resolve cada tiro por RK4 (com $z$ representando $y'$ no campo da função) e faz
sete iterações da secante. O controle $kk$ escolhe quantos tiros aparecem no gráfico.
\begin{enumerate}[nosep]
  \item Confira a tabela do exemplo e o erro da solução final em relação à exata.
  \item Mude o chute $s_0$ para $-20$ e para $10$. A secante ainda converge? Em quantas iterações?
  \item Aumente $n$ para 40. O $s$ final fica mais perto de $-14$? Por quê?
  \item PVC linear: digite $g=4(y-x)$, $a=0$, $b=1$, $y(a)=0$, $y(b)=2$, a solução exata
        $y=x+\operatorname{senh}(2x)/\operatorname{senh}(2)$ (digite \texttt{x + sinh(2x)/sinh(2)}) e
        $s_0=0$. Confirme que a secante acerta já em $s_2$.
\end{enumerate}
\end{atividade}

% ==================================================================
\section{Método das diferenças finitas}

\subsection{Aproximação das derivadas}

Dividimos $[a,b]$ em $N+1$ subintervalos de comprimento $h=\frac{b-a}{N+1}$, com nós
$x_i=a+ih$, $i=0,1,\dots,N+1$; os nós $x_1,\dots,x_N$ são os \emph{pontos internos}. Pela
fórmula de Taylor,
\[
y(x_{i\pm1})=y(x_i)\pm h\,y'(x_i)+\frac{h^2}{2}y''(x_i)\pm\frac{h^3}{6}y'''(x_i)+\frac{h^4}{24}y^{(4)}(\xi_\pm).
\]
Somando e subtraindo as duas expansões, obtemos as \emph{diferenças centradas}
\[
y''(x_i)=\frac{y(x_{i+1})-2y(x_i)+y(x_{i-1})}{h^2}+O(h^2),\qquad
y'(x_i)=\frac{y(x_{i+1})-y(x_{i-1})}{2h}+O(h^2).
\]

\subsection{O sistema linear}

Para o PVC linear $y''=p(x)y'+q(x)y+r(x)$, $y(a)=\alpha$, $y(b)=\beta$, trocamos as derivadas
pelas diferenças em cada ponto interno e $y(x_i)$ pela incógnita $w_i$
($w_0=\alpha$, $w_{N+1}=\beta$):
\[
\frac{w_{i+1}-2w_i+w_{i-1}}{h^2}=p_i\,\frac{w_{i+1}-w_{i-1}}{2h}+q_i\,w_i+r_i,
\]
com $p_i=p(x_i)$ etc. Multiplicando por $-h^2$ e agrupando:
\[
\boxed{\;\Bigl(-1-\frac h2p_i\Bigr)w_{i-1}+\bigl(2+h^2q_i\bigr)w_i+\Bigl(-1+\frac h2p_i\Bigr)w_{i+1}=-h^2r_i,\qquad i=1,\dots,N.\;}
\]
Na primeira e na última equação, $w_0=\alpha$ e $w_{N+1}=\beta$ são conhecidos e passam para o
segundo membro. O resultado é um sistema $N\times N$ \emph{tridiagonal}, que pode ser resolvido
por eliminação de Gauss sem trocas de linhas e com custo proporcional a $N$ (não a $N^3$).

\begin{observacao}
Se $q(x)\ge0$ e $h\,|p(x)|\le2$, a matriz é diagonalmente dominante: o sistema tem solução
única e os métodos de Gauss-Jacobi e Gauss-Seidel do Capítulo~\ref{cap:sistemas} convergem.
O erro do método é $O(h^2)$, como o das diferenças centradas: dobrar o número de
subintervalos divide o erro por 4.
\end{observacao}

\begin{exemplo}[: $y''=4(y-x)$, $y(0)=0$, $y(1)=2$]
Aqui $p=0$, $q=4$ e $r=-4x$; a solução exata é $y=x+\dfrac{\operatorname{senh}(2x)}{\operatorname{senh}2}$.
Com $N=3$ pontos internos, $h=0{,}25$ e $2+h^2q=2{,}25$; os segundos membros são
$-h^2r_i=0{,}25\,x_i$, isto é, $0{,}0625$; $0{,}125$; $0{,}1875$, e a última equação recebe
$+\beta=2$:
\[
\begin{pmatrix}2{,}25&-1&0\\-1&2{,}25&-1\\0&-1&2{,}25\end{pmatrix}
\begin{pmatrix}w_1\\w_2\\w_3\end{pmatrix}=
\begin{pmatrix}0{,}0625\\0{,}125\\2{,}1875\end{pmatrix}.
\]
A solução é $w_1=0{,}3951247$, $w_2=0{,}8265306$, $w_3=1{,}3395692$; os valores exatos são
$0{,}3936767$, $0{,}8240271$ e $1{,}3370861$, com erro máximo $2{,}5\times10^{-3}$. Aumentando $N$:
\begin{center}
\renewcommand{\arraystretch}{1.25}
\begin{tabular}{c|c|c}
$N$ & $h$ & erro máximo\\ \hline
3 & $0{,}25$ & $2{,}50\times10^{-3}$\\
4 & $0{,}2$ & $1{,}73\times10^{-3}$\\
9 & $0{,}1$ & $4{,}39\times10^{-4}$\\
19 & $0{,}05$ & $1{,}10\times10^{-4}$
\end{tabular}
\end{center}
De $h=0{,}1$ para $h=0{,}05$, o erro cai pelo fator $4$, como previsto.
\end{exemplo}

\begin{atividade}{diferenças finitas para PVC linear}{\arqDifFin}
O aplicativo monta e resolve o sistema tridiagonal. À direita aparecem a matriz aumentada
(para $N\le8$), a tabela com $x_i$, $w_i$, a solução exata e o erro, e o erro máximo.
\begin{enumerate}[nosep]
  \item Com $N=3$, confira a matriz e a solução do exemplo.
  \item Faça $N=4,9,19$ e confira a tabela de erros.
  \item Digite $p(x)=-2/x$, $q(x)=2/x^2$, $r(x)=\operatorname{sen}(\ln x)/x^2$, $a=1$, $b=2$,
        $y(a)=1$, $y(b)=2$ (desmarque \emph{Mostrar a solução exata}, que se refere ao exemplo
        anterior). Como as entradas da matriz mudam de uma linha para outra?
  \item Compare com o método do tiro aplicado ao mesmo PVC linear (atividade anterior, item 4).
        Qual método dá menor erro com $h=0{,}1$?
\end{enumerate}
\end{atividade}

% ==================================================================
\section{Comparação dos métodos}

\begin{center}
\renewcommand{\arraystretch}{1.4}
\begin{tabular}{l|c|p{3.4cm}|p{5.2cm}}
método & ordem & custo & observações\\ \hline
Euler & 1 & 1 avaliação de $f$ por passo & simples; base para entender os outros; exige $h$ pequeno\\
Heun / ponto médio & 2 & 2 avaliações por passo & previsão e correção\\
RK4 & 4 & 4 avaliações por passo & o método de uso geral mais popular\\
tiro (com RK4) & 4 & um PVI por iteração & reaproveita métodos de PVI; pode divergir\\
diferenças finitas & 2 & um sistema tridiagonal & robusto para PVC lineares
\end{tabular}
\end{center}

\paragraph{Em resumo.}
\begin{itemize}[nosep]
  \item Para \textbf{PVI}, os métodos de Runge-Kutta são de passo único (só usam $y_i$ para
        calcular $y_{i+1}$), fáceis de programar e de mudar de passo; o RK4 é o padrão em
        problemas não rígidos. Programas profissionais usam pares de Runge-Kutta de ordens
        diferentes para estimar o erro e ajustar $h$ automaticamente (passo adaptativo).
  \item A \textbf{estabilidade} limita o passo, independentemente da precisão desejada. Para
        problemas rígidos são necessários métodos implícitos.
  \item Para \textbf{PVC}, o método do tiro aproveita toda a precisão do RK4, mas pode ser
        sensível ao chute inicial e ao crescimento das soluções; as diferenças finitas levam
        a sistemas lineares bem comportados e se estendem naturalmente às equações
        diferenciais parciais.
  \item Tudo depende dos capítulos anteriores: Taylor e erros (Capítulo~\ref{cap:pontoflutuante}),
        secante (Capítulo~\ref{cap:zeros}), sistemas tridiagonais (Capítulo~\ref{cap:sistemas})
        e regras de integração (Capítulo~\ref{cap:integracao}).
\end{itemize}

% ==================================================================
\section{Exercícios}

\begin{enumerate}[label=\textbf{\arabic*.},leftmargin=*]
\item Aplique o método de Euler com $h=0{,}1$ ao PVI $y'=-2xy$, $y(0)=1$, até $x=0{,}5$, e
  compare com a solução exata $y=e^{-x^2}$.

\item Para o mesmo PVI, faça dois passos ($h=0{,}1$) pelo método de Heun e pelo RK4.

\item Mostre que, quando $f(x,y)=g(x)$ não depende de $y$, um passo do RK4 é a regra $1/3$ de
  Simpson e um passo de Heun é a regra dos trapézios no intervalo $[x_i,x_{i+1}]$.

\item Qual é o maior passo $h$ para o qual o método de Euler aplicado a $y'=-50y$, $y(0)=1$,
  produz aproximações que decaem para zero?

\item Escreva como sistema de primeira ordem: (a) $y''+4y'+3y=0$; (b) o pêndulo amortecido
  $\theta''+0{,}5\,\theta'+9{,}8\operatorname{sen}\theta=0$; (c) $y'''=x\,y-y'$.

\item Faça um passo do RK4 com $h=0{,}1$ para o oscilador $u'=v$, $v'=-u$, $u(0)=1$, $v(0)=0$,
  e compare com $(\cos0{,}1;\ -\operatorname{sen}0{,}1)$.

\item Verifique que $y=(x+1)^2-\tfrac12e^x$ é solução do PVI $y'=y-x^2+1$, $y(0)=0{,}5$.

\item Use diferenças finitas com $N=1$ (um único ponto interno) para o PVC $y''=4(y-x)$,
  $y(0)=0$, $y(1)=2$, e compare $w_1$ com $y(0{,}5)$.

\item Para o PVC linear $y''=4(y-x)$, $y(0)=0$, $y(1)=2$, determine o valor exato de $y'(0)$ e
  explique por que o método do tiro com a secante o encontra (a menos do erro do RK4) já em $s_2$.

\item \textbf{(GeoGebra)} No aplicativo de Euler, use $f(x,y)=-20y$, $y_0=1$, $n=10$ e
  $h=0{,}1$. Explique o resultado usando o fator $1+h\lambda$.

\item \textbf{(GeoGebra)} No aplicativo de sistemas, use o modelo de Lotka-Volterra com Euler e
  $h=0{,}1$, $n=200$ e depois $h=0{,}01$, $n=2000$. O ponto final de Euler se aproxima do RK4?
  Com que ordem?

\item \textbf{(GeoGebra)} No aplicativo de diferenças finitas, faça $N=4$, $9$ e $19$ e calcule
  as razões entre os erros máximos. Que ordem elas indicam?
\end{enumerate}

\subsection*{Respostas}
\begin{enumerate}[label=\textbf{\arabic*.},leftmargin=*,nosep]
\item $y_1=1$; $y_2=0{,}98$; $y_3=0{,}9408$; $y_4=0{,}884352$; $y_5=0{,}8136038$; exato
  $e^{-0{,}25}=0{,}7788008$ (erro $0{,}035$).
\item Heun: $y_1=0{,}99$, $y_2=0{,}960696$. RK4: $y_1=0{,}9900498$ (erro $4\times10^{-10}$),
  $y_2=0{,}9607894$ (erro $4\times10^{-9}$); exatos $0{,}9900498$ e $0{,}9607894$.
\item Com $f=g(x)$: $k_1=g(x_i)$, $k_2=k_3=g(x_i+\frac h2)$, $k_4=g(x_i+h)$, logo
  $y_{i+1}-y_i=\frac h6[g(x_i)+4g(x_i+\frac h2)+g(x_{i+1})]$, a regra $1/3$ com espaçamento
  $\frac h2$. Em Heun, $k_1=g(x_i)$ e $k_2=g(x_i+h)$, logo $\frac h2(k_1+k_2)$ é o trapézio.
\item $|1-50h|<1\iff h<0{,}04$.
\item (a) $u'=v$, $v'=-4v-3u$. (b) $u'=v$, $v'=-0{,}5v-9{,}8\operatorname{sen}u$.
  (c) $u'=v$, $v'=w$, $w'=xu-v$.
\item $(0{,}99500417;\ -0{,}09983333)$; exato $(0{,}99500417;\ -0{,}09983342)$.
\item $y'=2(x+1)-\frac12e^x$ e $y-x^2+1=2x+2-\frac12e^x$; além disso, $y(0)=1-\frac12=0{,}5$.
\item $h=0{,}5$: $(2+0{,}25\cdot4)\,w_1=0{,}25\cdot0{,}5\cdot4+0+2$, isto é, $3w_1=2{,}5$, logo
  $w_1=0{,}8333$; o exato é $0{,}8240$ (erro $0{,}0093$).
\item $y'(0)=1+\dfrac{2}{\operatorname{senh}2}\approx1{,}5514411$. A EDO é linear, logo
  $\phi(s)$ é afim e a secante dá o zero exato de $\phi$ no primeiro passo.
\item $1+h\lambda=1-20\cdot0{,}1=-1$: as aproximações alternam entre $-1$ e $1$ e não decaem
  (limite da estabilidade, $h=2/|\lambda|$).
\item Com $h=0{,}1$, Euler termina em $(3{,}32;\ 0{,}50)$; com $h=0{,}01$, em $(9{,}74;\ 1{,}20)$,
  mais perto do RK4 $(10{,}21;\ 1{,}33)$: a distância ao ponto do RK4 cai de cerca de $6{,}9$
  para $0{,}5$. Mesmo com 10 vezes mais passos, Euler continua bem menos preciso que o RK4
  com $h=0{,}1$.
\item As razões ficam perto de $1{,}73\times10^{-3}/4{,}39\times10^{-4}\approx3{,}9$ e
  $4{,}39\times10^{-4}/1{,}10\times10^{-4}\approx4{,}0$ (o passo cai à metade): ordem 2.
\end{enumerate}
